\subsection{MV/LV network transformers}
\label{sec:de-network}

Figure~\ref{fig:de-network} shows the vector groups of the network
transformer. The clock between the two diagrams gives the vector-group hour
of Equation~\eqref{eq:de-clock}.

\begin{figure}[!ht]
  \centering
  % generated by north-american-practice/fig/connection_figures.py -- do not edit
\begin{tikzpicture}[x=0.27mm, y=-0.27mm, line cap=round, line join=round]
  \draw[hv] (90.0,40.0) -- (142.0,130.0);
  \draw[hv] (142.0,130.0) -- (38.0,130.0);
  \node[font=\footnotesize, text=atInfo, anchor=center, inner sep=0.5pt] at (90.0,145.0) {20\,kV};
  \draw[hv] (38.0,130.0) -- (90.0,40.0);
  \fill[atInfo] (90.0,40.0) circle (3.0);
  \node[font=\footnotesize, text=atInfo, anchor=center, inner sep=0.5pt] at (90.0,27.0) {1U};
  \fill[atInfo] (142.0,130.0) circle (3.0);
  \node[font=\footnotesize, text=atInfo, anchor=center, inner sep=0.5pt] at (153.2,136.5) {1V};
  \fill[atInfo] (38.0,130.0) circle (3.0);
  \node[font=\footnotesize, text=atInfo, anchor=center, inner sep=0.5pt] at (26.8,136.5) {1W};
  \draw[lv] (318.0,100.0) -- (348.0,152.0);
  \draw[lv] (318.0,100.0) -- (258.0,100.0);
  \draw[lv] (318.0,100.0) -- (348.0,48.0);
  \node[font=\footnotesize, text=atAccentFg, anchor=center, inner sep=0.5pt] at (344.3,119.5) {230};
  \fill[atFg] (318.0,100.0) circle (3.0);
  \draw[gnd] (318.0,103.0) -- (318.0,113.0);
  \draw[gnd] (311.0,113.0) -- (325.0,113.0);
  \draw[gnd] (313.5,117.0) -- (322.5,117.0);
  \draw[gnd] (316.0,121.0) -- (320.0,121.0);
  \node[font=\footnotesize, text=atFg, anchor=east, inner sep=0.5pt] at (306.0,120.0) {2N};
  \fill[atAccentFg] (348.0,152.0) circle (3.0);
  \node[font=\footnotesize, text=atAccentFg, anchor=center, inner sep=0.5pt] at (354.5,163.2) {2U};
  \fill[atAccentFg] (258.0,100.0) circle (3.0);
  \node[font=\footnotesize, text=atAccentFg, anchor=center, inner sep=0.5pt] at (245.0,100.0) {2V};
  \fill[atAccentFg] (348.0,48.0) circle (3.0);
  \node[font=\footnotesize, text=atAccentFg, anchor=center, inner sep=0.5pt] at (354.5,36.8) {2W};
  \draw[aux, solid, fill=white] (188.0,58.0) circle (15);
  \draw[aux, solid] (188.0,43.0) -- (188.0,46.0);
  \draw[aux, solid] (195.5,45.0) -- (194.0,47.6);
  \draw[aux, solid] (201.0,50.5) -- (198.4,52.0);
  \draw[aux, solid] (203.0,58.0) -- (200.0,58.0);
  \draw[aux, solid] (201.0,65.5) -- (198.4,64.0);
  \draw[aux, solid] (195.5,71.0) -- (194.0,68.4);
  \draw[aux, solid] (188.0,73.0) -- (188.0,70.0);
  \draw[aux, solid] (180.5,71.0) -- (182.0,68.4);
  \draw[aux, solid] (175.0,65.5) -- (177.6,64.0);
  \draw[aux, solid] (173.0,58.0) -- (176.0,58.0);
  \draw[aux, solid] (175.0,50.5) -- (177.6,52.0);
  \draw[aux, solid] (180.5,45.0) -- (182.0,47.6);
  \draw[hv, line width=0.9pt] (188.0,58.0) -- (188.0,46.0);
  \draw[lv, line width=0.9pt] (188.0,58.0) -- (192.8,66.2);
  \draw[arr] (172.0,104.0) -- (204.0,104.0);
  \node[font=\footnotesize, text=atFgMuted, anchor=center, inner sep=0.5pt] at (188.0,122.0) {150\textdegree};
  \node[font=\footnotesize, text=atFgMuted, anchor=center, inner sep=0.5pt] at (90.0,196.0) {20\,kV $\Delta$};
  \node[font=\footnotesize, text=atFgMuted, anchor=center, inner sep=0.5pt] at (318.0,196.0) {400/230\,V, 2N earthed};
  \node[font=\small\bfseries, text=atFg, anchor=west, inner sep=0.5pt] at (-6.0,18.0) {(a) Dyn5};
  \draw[hv] (90.0,315.0) -- (90.0,255.0);
  \draw[hv] (90.0,315.0) -- (142.0,345.0);
  \draw[hv] (90.0,315.0) -- (38.0,345.0);
  \fill[atInfo] (90.0,255.0) circle (3.0);
  \node[font=\footnotesize, text=atInfo, anchor=center, inner sep=0.5pt] at (90.0,242.0) {1U};
  \fill[atInfo] (142.0,345.0) circle (3.0);
  \node[font=\footnotesize, text=atInfo, anchor=center, inner sep=0.5pt] at (153.2,351.5) {1V};
  \fill[atInfo] (38.0,345.0) circle (3.0);
  \node[font=\footnotesize, text=atInfo, anchor=center, inner sep=0.5pt] at (26.8,351.5) {1W};
  \draw[lv] (318.0,315.0) -- (318.0,349.6) -- (348.0,367.0);
  \draw[lv] (318.0,315.0) -- (288.0,297.7) -- (258.0,315.0);
  \draw[lv] (318.0,315.0) -- (348.0,297.7) -- (348.0,263.0);
  \fill[atFg] (318.0,315.0) circle (3.0);
  \draw[gnd] (318.0,318.0) -- (318.0,328.0);
  \draw[gnd] (311.0,328.0) -- (325.0,328.0);
  \draw[gnd] (313.5,332.0) -- (322.5,332.0);
  \draw[gnd] (316.0,336.0) -- (320.0,336.0);
  \node[font=\footnotesize, text=atFg, anchor=west, inner sep=0.5pt] at (330.0,335.0) {2N};
  \fill[atAccentFg] (348.0,367.0) circle (3.0);
  \node[font=\footnotesize, text=atAccentFg, anchor=center, inner sep=0.5pt] at (354.5,378.2) {2U};
  \fill[atAccentFg] (258.0,315.0) circle (3.0);
  \node[font=\footnotesize, text=atAccentFg, anchor=center, inner sep=0.5pt] at (245.0,315.0) {2V};
  \fill[atAccentFg] (348.0,263.0) circle (3.0);
  \node[font=\footnotesize, text=atAccentFg, anchor=center, inner sep=0.5pt] at (354.5,251.8) {2W};
  \draw[aux, solid, fill=white] (188.0,273.0) circle (15);
  \draw[aux, solid] (188.0,258.0) -- (188.0,261.0);
  \draw[aux, solid] (195.5,260.0) -- (194.0,262.6);
  \draw[aux, solid] (201.0,265.5) -- (198.4,267.0);
  \draw[aux, solid] (203.0,273.0) -- (200.0,273.0);
  \draw[aux, solid] (201.0,280.5) -- (198.4,279.0);
  \draw[aux, solid] (195.5,286.0) -- (194.0,283.4);
  \draw[aux, solid] (188.0,288.0) -- (188.0,285.0);
  \draw[aux, solid] (180.5,286.0) -- (182.0,283.4);
  \draw[aux, solid] (175.0,280.5) -- (177.6,279.0);
  \draw[aux, solid] (173.0,273.0) -- (176.0,273.0);
  \draw[aux, solid] (175.0,265.5) -- (177.6,267.0);
  \draw[aux, solid] (180.5,260.0) -- (182.0,262.6);
  \draw[hv, line width=0.9pt] (188.0,273.0) -- (188.0,261.0);
  \draw[lv, line width=0.9pt] (188.0,273.0) -- (192.8,281.2);
  \draw[arr] (172.0,319.0) -- (204.0,319.0);
  \node[font=\footnotesize, text=atFgMuted, anchor=center, inner sep=0.5pt] at (188.0,337.0) {150\textdegree};
  \node[font=\footnotesize, text=atFgMuted, anchor=center, inner sep=0.5pt] at (90.0,411.0) {20\,kV Y, no neutral};
  \node[font=\footnotesize, text=atFgMuted, anchor=center, inner sep=0.5pt] at (318.0,411.0) {400/230\,V zigzag};
  \node[font=\small\bfseries, text=atFg, anchor=west, inner sep=0.5pt] at (-6.0,233.0) {(b) Yzn5};
  \draw[hv] (90.0,470.0) -- (142.0,560.0);
  \draw[hv] (142.0,560.0) -- (38.0,560.0);
  \node[font=\footnotesize, text=atInfo, anchor=center, inner sep=0.5pt] at (90.0,575.0) {20\,kV};
  \draw[hv] (38.0,560.0) -- (90.0,470.0);
  \fill[atInfo] (90.0,470.0) circle (3.0);
  \node[font=\footnotesize, text=atInfo, anchor=center, inner sep=0.5pt] at (90.0,457.0) {1U};
  \fill[atInfo] (142.0,560.0) circle (3.0);
  \node[font=\footnotesize, text=atInfo, anchor=center, inner sep=0.5pt] at (153.2,566.5) {1V};
  \fill[atInfo] (38.0,560.0) circle (3.0);
  \node[font=\footnotesize, text=atInfo, anchor=center, inner sep=0.5pt] at (26.8,566.5) {1W};
  \draw[lv] (318.0,530.0) -- (288.0,478.0);
  \draw[lv] (318.0,530.0) -- (378.0,530.0);
  \draw[lv] (318.0,530.0) -- (288.0,582.0);
  \node[font=\footnotesize, text=atAccentFg, anchor=center, inner sep=0.5pt] at (291.7,510.5) {230};
  \fill[atFg] (318.0,530.0) circle (3.0);
  \draw[gnd] (318.0,533.0) -- (318.0,543.0);
  \draw[gnd] (311.0,543.0) -- (325.0,543.0);
  \draw[gnd] (313.5,547.0) -- (322.5,547.0);
  \draw[gnd] (316.0,551.0) -- (320.0,551.0);
  \node[font=\footnotesize, text=atFg, anchor=west, inner sep=0.5pt] at (330.0,550.0) {2N};
  \fill[atAccentFg] (288.0,478.0) circle (3.0);
  \node[font=\footnotesize, text=atAccentFg, anchor=center, inner sep=0.5pt] at (281.5,466.8) {2U};
  \fill[atAccentFg] (378.0,530.0) circle (3.0);
  \node[font=\footnotesize, text=atAccentFg, anchor=center, inner sep=0.5pt] at (391.0,530.0) {2V};
  \fill[atAccentFg] (288.0,582.0) circle (3.0);
  \node[font=\footnotesize, text=atAccentFg, anchor=center, inner sep=0.5pt] at (281.5,593.2) {2W};
  \draw[aux, solid, fill=white] (188.0,488.0) circle (15);
  \draw[aux, solid] (188.0,473.0) -- (188.0,476.0);
  \draw[aux, solid] (195.5,475.0) -- (194.0,477.6);
  \draw[aux, solid] (201.0,480.5) -- (198.4,482.0);
  \draw[aux, solid] (203.0,488.0) -- (200.0,488.0);
  \draw[aux, solid] (201.0,495.5) -- (198.4,494.0);
  \draw[aux, solid] (195.5,501.0) -- (194.0,498.4);
  \draw[aux, solid] (188.0,503.0) -- (188.0,500.0);
  \draw[aux, solid] (180.5,501.0) -- (182.0,498.4);
  \draw[aux, solid] (175.0,495.5) -- (177.6,494.0);
  \draw[aux, solid] (173.0,488.0) -- (176.0,488.0);
  \draw[aux, solid] (175.0,480.5) -- (177.6,482.0);
  \draw[aux, solid] (180.5,475.0) -- (182.0,477.6);
  \draw[hv, line width=0.9pt] (188.0,488.0) -- (188.0,476.0);
  \draw[lv, line width=0.9pt] (188.0,488.0) -- (183.2,479.8);
  \draw[arr] (172.0,534.0) -- (204.0,534.0);
  \node[font=\footnotesize, text=atFgMuted, anchor=center, inner sep=0.5pt] at (188.0,552.0) {330\textdegree};
  \node[font=\footnotesize, text=atFgMuted, anchor=center, inner sep=0.5pt] at (90.0,626.0) {20\,kV $\Delta$};
  \node[font=\footnotesize, text=atFgMuted, anchor=center, inner sep=0.5pt] at (318.0,626.0) {400/230\,V, 2N earthed};
  \node[font=\small\bfseries, text=atFg, anchor=west, inner sep=0.5pt] at (-6.0,448.0) {(c) Dyn11};
\end{tikzpicture}

  \caption{Vector groups of MV/LV network transformers 20\,kV --
  400/230\,V: (a) Dyn5, (b) Yzn5, (c) Dyn11. Clock: higher-voltage phasor
  of phase U at 12, lower-voltage phasor at the clock hour.}
  \label{fig:de-network}
\end{figure}

\subsubsection{Dyn5}
Dyn5 is the vector group of German MV/LV network transformers,
Figure~\ref{fig:de-network}(a). The wye secondary is solidly earthed at 2N
and supplies 230\,V from line to neutral. The delta primary provides a
closed path for the zero-sequence current of unbalanced load and of earth
faults on the low-voltage side, so that the zero-sequence impedance seen
from the low-voltage side is low and no zero-sequence current flows in the
medium-voltage network. Network transformers up to 36\,kV are specified to
EN~50588-1~[15], and their losses are limited by the Ecodesign
regulation~[16].

\subsubsection{Yzn5}
In the zigzag secondary of Figure~\ref{fig:de-network}(b), each phase
consists of two half-windings on two different limbs. A single-phase load
then magnetises two limbs, and its ampere-turns are balanced on the primary
side without a delta winding. Yzn5 is used for small network transformers
that supply predominantly single-phase load.

\subsubsection{Dyn11}
Dyn11 has the same windings as Dyn5, connected for a lower-voltage side
leading by 30\,$^{\circ}$, Figure~\ref{fig:de-network}(c). It is the usual
vector group in a number of other countries and occurs in German industrial
installations. A Dyn5 and a Dyn11 transformer are operated in parallel only
after two phases are interchanged on both sides of one of them~[17].
\FloatBarrier
